\newcommand{\paths}{P^*}

\section{Algorithms for the sparsest $K$-bounded cut}
\label{sec:alpha}

Since solving tasks related to sparsest cut is easier on trees, an efficient approach to solving such problems is to first embed the graph into a tree while partially preserving the value of all cuts. To do so, one can use the cut-sparsifier tree decomposition of Racke \cite{OptimalHierarchicalDecompositionsForCongestionMinimizationInNetworks}. This decomposition allows for an embedding of any graph into a probability distribution over a polynomial number of trees such that for every cut, the expected value of the cut in the tree is at most $O\br{\log n}$ times the value of the cut in the original graph, and for every cut, its value in the original graph is at most its value in every tree in the support of the distribution. Therefore, if we can solve our problem on trees, then we can solve it on general graphs with an $O\br{\log n}$-approximation ratio by first embedding it into a distribution over trees and then solving it on each tree in the support of the distribution. Then, the approximation ratio is achieved by choosing the best cut among all cuts produced for these trees. In what follows we make this ideas more precise.

Lemmas~\ref{lem:mu} and~\ref{lem:single-edge} proven below give the following theorem that is the main result of this section.
\begin{theorem} \label{thm:phiK}
  There exists an $\alpha=\bigo\br{\log n}$-approximation algorithm for general graphs and an exact algorithm for trees for the sparsest $K$-bounded cut problem, for any choice of $K$.
\end{theorem}
Our main tool for proving the result is the notion of cut sparcifiers.
We recall the necessary notation in the next section.

\subsection{Decomposition trees and multicommodity flows}
Given a graph $G=(V,E,c)$, denote by $\paths$ the set of all paths in $G$.
A \emph{decomposition tree} for $G$ is a rooted tree
$T=(V_T,E_T, c_T)$ together with:
\begin{enumerate}
  \item A node map $m_V:V_T\to V$ that is bijective on the leaves of $T$ and $V$.
  \item An edge map $m_E:E_T\to \paths$ assigning each tree edge $uv$ to a path in $G$ between the mapped endpoints $m_V(u)$ and $m_V(v)$.
  \item Induced maps $m'_V:V\to V_T$ which maps each vertex of $V$ to the corresponding leaf of $T$ and $m'_E:E\to \paths_T$ which maps an edge $uv\in E$ to the unique (shortest) path between $m'_V(u)$ and $m'_V(v)$ in $T$.
  \item For a tree edge $e=uv\in E_T$, let $V_{u},V_{v}$ be the leaf partition induced by deleting $e$ (i.e., a leaf $y$ is in $V_u$ (respectively $V_v$) if and only if the endpoint of $e$ that is closer to $y$ in $T$ is $u$ (respectively $v$)).
  The \emph{capacity} of $e$ is then defined as
\[
  c_T(e):=\sum_{x\in V_{u}, y\in V_{v}} c(xy).
\]
\end{enumerate}
\DD{uwaga: w powyzszym zmienilem $c(x,y)$ na $c(xy)$ by parametrem byla krawedz jak wczesniej.}

A \emph{multicommodity flow} $f$ in a graph $G$ is a collection of flows $f^{(i)}:E\to\mathbb{R}_{\ge 0}$, $i\in\brc{1,\ldots,k}$, where each flow $f^{(i)}$ routes $d_i$ units between a source-sink pair $(s_i,t_i)$. The flows needs to satisfy the demand constraints and should respect edge capacities.
\DD{dobrze byloby dodefiniowac co to oznacza (a moze usunac ostatnie zdanie?)}

The \emph{congestion} of a multicommodity flow $f$ is a way to measure how much these capacities are violated.
Formally, we define it as the maximum ratio of the total flow to capacity over all edges:
\[
  \xi_G\br{f} := \max_{e\in E} \frac{\sum_{i=1}^{k} f^{(i)}(e)}{c(e)}.
\]
A multicommodity flow is \emph{feasible} if its congestion is at most $1$. Given a flow $f$ in $G$, we can map it to a flow in a decomposition tree $T$ using the edge map $m_E$.
The \emph{mapped flow} $m'(f)$ for each tree edge $e_t\in E_T$ sums the flow of all edges in $G$ that are mapped to paths in $T$ that include $e_t$. Formally:
\[
  m'(f)(e_t) := \sum_{e\in E: e_t\in m'_E(e)} f(e).
\]

R\"{a}cke's theorems relate the congestion of flows in a graph to the congestion in decomposition trees:

\begin{theorem}[R\"{a}cke (STOC 2008, Theorem 2)]
\label{thm:Racke2}
Let $f$ be any multicommodity flow in $G$ with congestion $\xi_G\br{f}$.
For any decomposition tree $T$ of $G$, the mapped flow $m'(f)$ in $T$ has congestion $\xi_T\br{m'\br{f}}\le \xi_G\br{f}$.
\end{theorem}

\begin{theorem}[R\"{a}cke (STOC 2008, Theorem 4)]
\label{thm:Racke4}
There exists a polynomial-time computable convex combination of decomposition
trees $(T_i,\lambda_i)$, $\sum_i\lambda_i=1$, such that for any family of multicommodity flows $f_i$ with congestion at most $1$ in $T_i$, the multicommodity flow in $G$ defined as
\[
  f:=\sum_i \lambda_i m(f_i)
\]
has congestion at most $\rho=O(\log n)$ in $G$.
\end{theorem}

\subsection{Algorithms for the sparsest $K$-bounded cut}

\DD{czy to sprzedaje cala intuicje?:}
We start by converting the above theorems regarding the congestion into a statement regarding the cuts, which is more useful for our purposes.
\begin{theorem}[Dual cut bound from R\"{a}cke]\label{thm:racke-dual}
Let $G=(V,E,c)$ and let $(T_i,\lambda_i)$ be as in Theorem~\ref{thm:Racke4}.
Define $\mu$ by $\Pr_{T\sim\mu}[T=T_i]=\lambda_i$ and set
$\rho=\bigo(\log n)$. Then for every cut $U\subsetneq V$:
\begin{enumerate}[label=\textup{(\roman*)}]
  \item $c\br{E\br{U,\bar{U}}}\le c_T\br{E_T\br{U,\bar{U}}}$ for every $T$ in the support of $\mu$.
  \item $\Etree\left[c_T\br{E_T\br{U,\bar{U}}}\right]\le \rho\cdot c\br{E\br{U,\bar{U}}}$.
\end{enumerate}
%The proof is deferred to Appendix~\ref{app:racke-proof}.
\end{theorem}
\begin{proof}
Fix a cut $U\subsetneq V$.

For (i), let $F^*=c(E\br{U,\bar{U}})$. By max-flow/min-cut in $G$, there is a
single-commodity flow from $U$ to $\bar U$ of value $F^*$ and congestion $1$.
By Theorem~\ref{thm:Racke2}, the corresponding mapped flow in any $T$ has congestion at most $1$.
Hence this flow does not exceed the capacity
of the tree cut $\br{U,\bar{U}}$. Therefore,
\[
  F^*=c\br{E\br{U,\bar{U}}}\le c_T\br{E_T\br{U,\bar{U}}}.
\]

For (ii), for each tree $T_i$ choose a maximum flow $f_i$ from $U$ to $\bar{U}$ in $T_i$.
Since $T_i$ is a tree with edge capacities $c_{T_i}$, max-flow/min-cut principle gives that the value of the flow $f_i$ does not exceed $c_{T_i}\br{E_{T_i}\br{U,\bar{U}}}$.
% \[
%   \mathrm{val}(f_i)=c_{T_i}\br{E\br{U,\bar{U}}},
%\]
Hence, $f_i$ has congestion at most $1$ in $T_i$.
By Theorem~\ref{thm:Racke4} applied to the entire family of flows $f_i$, the mapped convex combination
\[
  f:=\sum_i \lambda_i m(f_i)
\]
has congestion at most $\rho$ in $G$. Its value is
\[
  F:=\sum_i \lambda_i\,\mathrm{val}(f_i)
   =\sum_i\lambda_i c_{T_i}(E\br{U,\bar{U}})
   =\Etree[c_T(E\br{U,\bar{U}})],
\]
where $\mathrm{val}(f_i)$ gives the value of a flow $f_i$.
Any flow from $U$ to $\bar{U}$ of value $F$ in $G$ must have congestion at least
$F/c(E\br{U,\bar{U}})$ (again by max-flow/min-cut), so
\[
  \frac{F}{c(E\br{U,\bar{U}})}\le \rho.
\]
Substituting $F=\Etree[c_T(E\br{U,\bar{U}})]$ gives
\[
  \Etree[c_T(E_T\br{U,\bar{U}})]\le \rho\,c(E\br{U,\bar{U}}).
\]
This proves (ii).
\end{proof}

\begin{lemma}\label{lem:mu}
  Let $G=(V,E,c,w)$ be any graph and let $\mu$ be the distribution from Theorem~\ref{thm:racke-dual}.
  there exists a tree $T$ in the support of $\mu$ such that $\OPTK(T) = \bigo\br{\log n}\cdot\OPTK(G)$.
\end{lemma}

\begin{proof}
  Let $\br{X^*,\bar{X}^*}$ be an optimal cut for generalized sparsest $K$-bounded cut.
  Since the demand function $w$ is the same in $G$ and in every tree $T$,
  our cut satisfies $0<w_{T}\br{X^*,\bar{X}^*}\le K$, hence the cut $X^*$ is feasible in every tree.
  Therefore by Theorem~\ref{thm:racke-dual},
  \begin{align*}
    \Etree\!\left[\OPTK(T)\right]
    \;&\le\;
    \Etree\!\left[\frac{c_T\br{E_T\br{X^*,\bar{X}^*}}}{w_T\br{X^*,\bar{X}^*}}\right]
    \;
    \\&=\;
    \frac{1}{w\br{X^*,\bar{X}^*}}\,\Etree\!\left[c_T\br{E_T\br{X^*,\bar{X}^*}}\right]
    \;\\&\le\;
    \frac{\rho\,c\br{E\br{X^*,\bar{X}^*}}}{w\br{X^*,\bar{X}^*}}
    \;\\&=\;
    \rho\cdot\OPTK(G).
  \end{align*}
  By averaging, some tree $T$ in the support satisfies
  $\OPTK(T)\le \rho\,\OPTK(G)$.
\end{proof}

%\subsection{Solving the \DD{sparsest cut} problem on a tree}
It remains to provide an algorithm for the sparsest $K$-bounded cut for trees.
It is a well known fact that an optimal sparsest cut in a tree is achieved by a single edge removal.
Below, we show that this is also the case for the $K$-bounded generalization.
\DD{dodalem def:}
For each edge $e$ of a tree $T$ we define $w_e$ to be the total demand cut by removing $e$ from $T$.
\begin{lemma}[Tree optimum is attained by a single edge]\label{lem:single-edge}
For a tree $T=(V_T,E_T,c_T,w_T)$,
  \[
    \OPTK(T)
    \;=\;
    \min_{\substack{e\in E_T \\ 0<w_e\le K}} \frac{c_T(e)}{w_e}.
  \]
\end{lemma}
\begin{proof}
Root $T$ at $r$.
For each edge $e$, let $A_e$ be the component away from $r$.
%  Define $c(e):=c_T(e)$ and $w_e:=w(\cut_T(A_e))$. Let $S$ be any feasible cut with $w(\cut(S))\le K$.
Take an arbitrary cut $(S,\bar{S})$ and denote for brevity $E'=E_T(S,\bar{S})$.
Then,
  \[
    w_T\br{S,\bar{S}} \leq \sum_{e\in E'} w_e.
  \]
  Since $w_T\br{S,\bar{S}}\le K$ and all $w_e\ge 0$, we have
  $w_e\le w_T\br{S,\bar{S}}\le K$ for every $e\in E'$.
  Also $w_e>0$ whenever $e\in E'$ (otherwise the edge can be trivially dropped from $E'$).
  Therefore,
  \[
    \frac{c_T(E')}{w_T\br{S,\bar{S}}}
    \geq
    \sum_{e\in E'} \frac{w_e}{\sum_{f\in E'}w_f}\cdot\frac{c_T(e)}{w_e}
    \;\ge\;
    \min_{e\in E_T, 0<w_e\le K} \brc{\frac{c_T(e)}{w_e}},
  \]
  where the inequality is due to the fact that the left-hand side is a weighted sum with coefficients $w_e/\sum_{f\in E'}w_f$ that sum up to $1$.
  In particular, we observe that by averaging, at least one of the terms in the sum must be at most the minimum of the terms.
\end{proof}

\begin{corollary} \label{cor:single-edge}
  For any tree $T$, $\OPTK(T)$ can be computed in $\bigo(n)$ time by a single scan over the edges of $T$.
\end{corollary}

\begin{proof}[Proof of Theorem~\ref{thm:phiK}]
By Theorem~\ref{thm:Racke4}, the convex combination of decomposition trees $\br{T_i,\lambda_i}$ can be computed in polynomial time.
Lemma~\ref{lem:mu} says that one of these trees ensures the approximation ratio of $\bigo\br{\log n}$ for the input graph $G$.
Moreover, Corollary~\ref{cor:single-edge} gives that the sparsest $K$-bounded cut can be solved optimally for all trees in the family.
\end{proof}


% \section{Semi-definite programming formulation of $h\br{S}$}
% In order to obtain an approximation algorithm for the problem of minimizing $h\br{S}=\frac{c\br{S, \bar{S}}}{w\br{S, V}}$, we will use a semi-definite programming relaxation. This problem can be formulated as the following quadratic program:

% \begin{align}
% \min \quad & \frac{1}{2K} \sum_{u,v \in V} c(u,v)\cdot(x_u-x_v) ^2\\[6pt]
% \text{s.t.} \quad 
% & x_u^2 + y_u^2 = 1 \quad \forall u \in V \\[6pt]
% & x_u\cdot y_u= 0 \quad \forall u \in V \\[6pt]
% & \sum_{u,v \in V} w(u,v)\cdot\, x_u^2 = K
% \end{align}

% Note, that in polynomial time, we can guess the value of $K=w(S^*, V)$ up to an arbitrary multiplicative precision.

% We now relax this QP to an SDP, allowing the variables to be vectors in unbounded (but polynomial sized) space. Furthermore, we add more constraints in order to strengthen the relaxation. We add triangle inequalities in order to ensure $x$ to be a metric. Moreover, we force all the vectors to be centered at $v/2$ where $v$ is the unit vector and we constraint all of them to be of length $1/2$. The SDP is as follows:
% \begin{align}
% \min \quad & \frac{1}{2K} \sum_{u,v \in E} c(u,v)\cdot(x_u-x_v) ^2 \\[6pt]
% \text{s.t.} \quad 
% & x_u^2 + y_u^2 = 1 \quad \forall u \in V \\[6pt]
% & x_u\cdot y_u= 0 \quad \forall u \in V \\[6pt]
% & v^2= 1 \\[6pt]
% & v\cdot x_u = x_u^2, v\cdot y_u = y_u^2 \quad \forall u \in V  \\[6pt]
% & \sum_{u,v \in V} w(u,v)\cdot\, x_u^2 = K \\[6pt]
% & (x_u - x_v)^2 \leq (x_u - x_w)^2 + (x_v - x_w)^2 \quad \forall u,v,w \in V
% \end{align}

% Define a metric on $V$ as $d(u,v)=(x_u-x_v)^2$. A map $f\colon X \to \mathbb{R}$ is called 1-Lipshitz if for all $x,y \in V$ $|f(x)-f(y)|\leq d(x,y)$. Given a 1-Lipschitz map $f$ and a capacity function $c$, define its average distortion under $c$ as:
% \[
% a_c(f) = \frac{\sum_{x,y \in X}c(x,y)\cdot d(x,y)}{\sum_{x,y \in X}c(x,y)\cdot |f(x)-f(y)|}
% \]

% Bourgains theorem states that for every such $w$ there exists an efficiently computable map $f$ such that $a_c(f)=\bigo(\log n)$. Observe, that the objective function of the SDP can be rewritten as $$\frac{1}{2K} \sum_{u,v \in V} c(u,v)\cdot d(u,v)$$

% Therefore, we know that there exists an efficiently computable map $f$ such that the following holds:
% $$
% \frac{1}{2K} \sum_{u,v \in V} c(u,v)\cdot d(u,v)\geq \frac{1}{2K} \sum_{u,v \in V} c(u,v)\cdot |f(x)-f(y)|\geq \frac{1}{a_c(f)\cdot 2K} \sum_{u,v \in V} c(u,v)\cdot d(u,v)
% $$
% where the first inequality is by the fact that $f$ is 1-Lipshitz and the second inequality is by definition of $a_c(f)$.

% Since by Bourgains theorem $f$ with $a_c(f)=\bigo(\log n)$ can be found in polynomial time we have that the objective function of the SDP embeds into $\mathbb{R}$ with $O(\log n)$ distortion.

% The remaining objective is to use the fact in design of a rounding procedure for the above SDP.

